% =====================================================================================================
%  Briefing report: Manoeuvre Pattern-of-Life Prediction for SDA -- everything this project has done
%  Compile:  pdflatex briefing.tex  (twice, for the contents and cross-references)
%  Figures are read from ../outputs/report/ (produced by `python main.py`). A missing figure is replaced
%  by a boxed note, so the document always compiles. On Overleaf, upload this file and the folder
%  outputs/report/ keeping the same relative layout (or change \graphicspath below).
% =====================================================================================================
\documentclass[11pt,a4paper]{article}

\usepackage[margin=2.3cm]{geometry}
\usepackage[T1]{fontenc}
\usepackage{lmodern}
\usepackage{amsmath,amssymb,bm}
\usepackage{graphicx}
\usepackage{booktabs,tabularx,longtable}
\usepackage{xcolor}
\usepackage{float}
\usepackage{caption}
\usepackage{enumitem}
\usepackage[colorlinks=true,linkcolor=blue!50!black,citecolor=blue!50!black,urlcolor=blue!50!black]{hyperref}

\graphicspath{{../outputs/report/}}
\setlength{\parskip}{4pt}
\setlength{\emergencystretch}{3em}   % lets lines with long code names break cleanly
\setlist{itemsep=2pt,topsep=3pt}

% ---- helpers -----------------------------------------------------------------------------------------
\newcommand{\code}[1]{\texttt{\detokenize{#1}}}
\newcommand{\figfile}[4]{% #1 file  #2 width  #3 caption  #4 label
  \begin{figure}[H]\centering
  \IfFileExists{../outputs/report/#1}{\includegraphics[width=#2\textwidth]{#1}}%
    {\fbox{\parbox{0.8\textwidth}{\centering\small Figure \code{#1} not found -- run \code{python main.py} first.}}}
  \caption{#3}\label{#4}\end{figure}}
\newsavebox{\briefbox}
\newenvironment{keybox}[1]%
  {\par\medskip\noindent\begin{lrbox}{\briefbox}\begin{minipage}{0.96\textwidth}\textbf{#1}\par\smallskip}%
  {\end{minipage}\end{lrbox}\fcolorbox{blue!40!black}{blue!4}{\usebox{\briefbox}}\par\medskip}
\newenvironment{whybox}%
  {\par\smallskip\noindent\begin{lrbox}{\briefbox}\begin{minipage}{0.96\textwidth}\textbf{Why this way?}\enspace}%
  {\end{minipage}\end{lrbox}\fcolorbox{orange!60!black}{orange!5}{\usebox{\briefbox}}\par\smallskip}

\title{\textbf{Manoeuvre Pattern-of-Life Prediction for Space Domain Awareness}\\[4pt]
\large A sequential briefing on every model, method and result of the project}
\author{UROP project (Imperial College London) --- briefing compiled from the code and its outputs}
\date{\today}

\begin{document}
\maketitle
\begin{abstract}
\noindent This document explains, in the order in which the work was done, how a Starlink-like satellite
that keeps its orbit inside a 500\,m semi-major-axis (SMA) deadband was simulated, how its manoeuvres were
detected from noisy tracking data, how the control law behind them was learned, and how that law was
used to forecast future manoeuvres. It summarises the physics (two-body gravity, the $J_2$ oblateness
term, atmospheric drag, lunar and solar third-body attraction, solar radiation pressure), the controller
(a smooth \emph{tanh} deadband with memory), the estimation methods (finite-difference regression, the
energy method, shooting, Kalman filtering, spectral analysis), the classification methods (Gaussian
mixture models and alternatives, scored per burn \emph{event}), the learning of the control law (SINDy and
a Bayesian variant) and nine extension studies aimed at the UK Space Agency brief (space weather,
realistic tracking passes, conjunction screening, sensor tasking, multiple control laws, fleets,
calibrated intervals, GEO and a library ``zoo'' for SINDy). Every number quoted comes from the latest run
of \code{python main.py} (\code{outputs/results.json}).
\end{abstract}

\tableofcontents
\newpage

% =====================================================================================================
\section{The brief and the question}
% =====================================================================================================
The UK Space Agency brief ``Manoeuvre Pattern of Life Prediction for SDA'' starts from a simple
observation: of the $\sim$50\,000 tracked objects about 15\,000 are active and may manoeuvre, and by 2035
mega-constellations will push this beyond 50\,000. Unpredicted manoeuvres break three core Space Domain
Awareness (SDA) functions:
\begin{enumerate}
  \item \textbf{close-approach (conjunction) screening}, which relies on predicted trajectories being
        accurate over the screening window --- in the brief's Figure~1 a Starlink satellite keeping its SMA
        inside a 500\,m deadband changes its miss distance by \emph{hundreds of km in a week};
  \item \textbf{track association}, because a manoeuvre makes the predicted position wrong, so new
        observations cannot be linked to the right catalogue object;
  \item \textbf{sensor scheduling}, because the uncertainty of an unpredicted manoeuvre has to be shrunk
        by tasking extra observations.
\end{enumerate}
Methods to detect \emph{past} manoeuvres exist; the gap the brief identifies is predicting \emph{future}
manoeuvre patterns by \textbf{reconstructing the analytical control law from time series of measurements}.

\begin{keybox}{The project in one sentence}
Simulate a station-keeping satellite with realistic physics, then show --- step by step and with honest
error bars --- that its burns can be detected from noisy data, that the control law behind them (``burn
when the mean SMA reaches the lower edge, stop at the upper edge'') can be learned, and that the learned
law forecasts future burns well enough to help screening, association and tasking.
\end{keybox}

\subsection{How the work progressed}
\begin{description}
  \item[Weeks 2--4] Basic simulation (gravity + drag + $J_2$), an \emph{impulsive} deadband controller,
        inverse fitting of the drag coefficient, an FFT test for the $J_2$ signature, simple burn detectors
        (threshold, CUSUM, $k$-means, GMM), a held-out comparison and a neural-network extrapolation test.
  \item[Week 5] A \emph{smooth} controller (tanh switches, a 7th state for the thruster, the mean SMA
        including the $J_2$ potential), burn detection in ``$\ddot r$ / $\dot v$ space'' with a GMM, and
        event-level scoring.
  \item[Report steps 1--9] A complete, documented pipeline on 400-day simulations: dynamics, regression,
        drag and noise, spectra and Kalman filtering, deadband control, classification, robustness,
        repetition (stacking and Gaussian processes) and SINDy/BINDy control-law learning with forecasts.
  \item[Step 0 and audits] Bias--variance tuning of every adjustable setting, walk-forward /
        walk-backward choice of the derivative method, a common figure style, and an audit of all figures.
  \item[Steps 10--18] Extensions towards the brief, using a fast mean-element model calibrated on the
        full simulation.
  \item[Step 19] Learning the thrust law and the natural dynamics together or apart, compared with the
        true model, and learning the natural law with no deadband.
\end{description}
All the old workstream stages were later re-written to use the \emph{same} 400-day data as the report, so
every figure in \code{outputs/report/} comes from one consistent data set.

% =====================================================================================================
\section{Units, reference orbit and numerical integration}
% =====================================================================================================
\subsection{Non-dimensional units}
The simulation is non-dimensional: the length unit is the Earth radius and the time unit is the inverse
Earth spin rate,
\begin{equation}
  \mathrm{DU} = 6371\ \text{km},\qquad \mathrm{TU} = 1/\omega_\oplus = 13\,713.44\ \text{s},\qquad
  \mu = \mu_\oplus\,\mathrm{TU}^2/\mathrm{DU}^3 = 289.873 .
\end{equation}
Working in these units keeps all state components of order one, which helps the ODE solvers and the
least-squares fits (Section~\ref{sec:regression} shows why badly scaled columns matter).

\subsection{Reference orbit}
A Starlink-like circular orbit at 550\,km altitude (SMA $\approx 6922$\,km), inclination $53^\circ$, period
$T\approx 95.5$\,min. The station-keeping band is $[6921.9,\,6922.4]$\,km in \emph{mean} SMA (500\,m wide),
the same as Figure~1 of the brief.

\subsection{Integrators}
The uncontrolled runs use \code{DOP853} (8th-order Runge--Kutta) with relative and absolute tolerance
$10^{-12}$ (non-dimensional). The controlled runs use \code{LSODA}, which switches automatically between a
non-stiff Adams method and a stiff BDF method: the tanh switches of the controller make the equations stiff
for a few seconds around every burn edge and non-stiff in between. A solver cross-check in step~1 confirms
that different integrators agree.

% =====================================================================================================
\section{The dynamical model (report step 1)}\label{sec:dynamics}
% =====================================================================================================
\subsection{Equations of motion}
The state is $\mathbf{x} = (\mathbf{r},\mathbf{v})$ (plus the thruster state $s$ for the controlled
runs). The acceleration is the sum of the forces:
\begin{equation}
  \dot{\mathbf r} = \mathbf v,\qquad
  \dot{\mathbf v} = \underbrace{-\frac{\mu\,\mathbf r}{r^3}}_{\text{two-body}}
   + \mathbf a_{J_2} + \mathbf a_{\text{drag}} + \mathbf a_{\text{Moon}} + \mathbf a_{\text{Sun}} + \mathbf a_{\text{SRP}}
   + \mathbf a_{\text{thrust}} .
\end{equation}

\paragraph{Earth oblateness ($J_2$).} The dominant perturbation at 550\,km comes from the equatorial
bulge. With $R_e$ the \textbf{equatorial} radius (6378.137\,km) and $J_2 = 1.08263\times10^{-3}$,
\begin{equation}
  \mathbf a_{J_2} = \frac{3}{2}\frac{J_2\,\mu R_e^2}{r^4}
  \begin{pmatrix} \tfrac{x}{r}\left(5\tfrac{z^2}{r^2}-1\right)\\[2pt]
                  \tfrac{y}{r}\left(5\tfrac{z^2}{r^2}-1\right)\\[2pt]
                  \tfrac{z}{r}\left(5\tfrac{z^2}{r^2}-3\right)\end{pmatrix},
  \qquad
  U_{J_2} = \frac{\mu J_2 R_e^2}{2 r^3}\left(3\sin^2\phi - 1\right),
\end{equation}
where $\phi$ is the latitude; $\mathbf a_{J_2} = -\nabla U_{J_2}$ (checked by a unit test).

\begin{keybox}{A bug found and fixed: the $J_2$ reference radius}
The code originally used the \emph{mean} Earth radius (6371\,km) in the $J_2$ term. $J_2$ is defined with
the \emph{equatorial} radius, so the perturbation was $(6371/6378.137)^2 \approx 0.22\%$ too weak. This
was fixed everywhere. The check (step~1): the theoretical nodal regression rate
\begin{equation}
  \dot\Omega = -\frac{3}{2}\,n\,J_2\left(\frac{R_e}{a}\right)^2\cos i
\end{equation}
integrated with the run's own (slowly drag-decaying) $a(t)$ agrees with the simulated right ascension of
the ascending node to \textbf{0.04\%} after 400 days ($-4.526^\circ$/day at the start, $-4.642^\circ$/day at the
end, the change being due to drag lowering the orbit).
\end{keybox}

\paragraph{Atmospheric drag.} With a constant (lumped) ballistic coefficient $c_d$,
\begin{equation}
  \mathbf a_{\text{drag}} = -c_d\,|\mathbf v|\,\mathbf v .
\end{equation}
Two values exist in the code: the Week~2--4 demonstration value \code{CD_TRUE} (about 45\,km/day of SMA
decay, which would deorbit the satellite in $\sim$12 days) and the \emph{realistic} value \code{SMOOTH_CD}
($\approx$0.13\,km/day, i.e.\ 5.3\,m/hour). Every 400-day run uses the realistic value. Steps~10--16 add a
time-varying density (space weather).

\paragraph{Third bodies (Moon, Sun).} The tidal (differential) attraction of a body at $\mathbf r_b$ with
gravitational parameter $\mu_b$ is
\begin{equation}
  \mathbf a_b = \mu_b\left(\frac{\mathbf r_b-\mathbf r}{|\mathbf r_b-\mathbf r|^3} - \frac{\mathbf r_b}{|\mathbf r_b|^3}\right),
\end{equation}
the second term removing the attraction of the Earth itself (the origin is Earth-centred). Moon and Sun move
on circular equatorial orbits (27.32 and 365.25 days).

\paragraph{Solar radiation pressure (SRP).} A flat plate always facing the Sun, no eclipses:
$|\mathbf a_{\text{SRP}}| = P\,C_R\,A/m$ with $P=4.56\,\mu$N/m$^2$, $C_R=1.3$, $A/m = 5/260$\,m$^2$/kg,
directed away from the Sun.

\subsection{The force budget}
\figfile{step1_dynamics/step1_summary.png}{0.98}{Step 1 summary: orbit, ground tracks, the $J_2$ nodal
regression check, the force budget, and how far each force moves the satellite over 15--360 days.}{fig:s1}
At 550\,km the mean accelerations are: gravity 8.3\,m/s$^2$, $J_2$ $1.2\times10^{-2}$, Moon
$8.8\times10^{-7}$, drag $8.1\times10^{-7}$, Sun $4.0\times10^{-7}$ and SRP $1.1\times10^{-7}$\,m/s$^2$.
$J_2$ is four orders of magnitude above everything else, and the thruster (about $2\times10^{-4}$\,m/s$^2$)
sits between $J_2$ and drag. Over a year the $J_2$-free orbit drifts by the orbit's diameter
($\sim 10^4$\,km) relative to the true one, while the Moon, Sun and SRP each move it by 10--30\,km.

\subsection{The 400-day data sets}
All later steps read the same cached runs (\code{deadband/long_run.py}, about 500\,MB in \code{.cache/}):
\begin{center}\small
\begin{tabularx}{\textwidth}{lXl}
\toprule
Run & Content & Used by\\\midrule
natural run & all forces, realistic drag, no thruster, 1-min samples & steps 1--4\\
divergence runs & gravity + drag, plus each extra force on its own & step 1\\
controlled run & smooth deadband controller, $J_2$ + drag, 6-s samples (102 burns) & steps 5--9\\
drag-only run & the controlled orbit without a thruster & step 5\\
\bottomrule
\end{tabularx}
\end{center}
Steps 6--9 learn from days 0--40 (TRAIN) and are tested on days 40--400 (TEST); results are reported at
horizons of 15, 30, 60, 120, 240 and 360 days.

% =====================================================================================================
\section{The station-keeping controller (report step 5)}\label{sec:controller}
% =====================================================================================================
\subsection{Why the Week 4 controller had to change}
The Week~4 controller fired an \emph{impulsive} raise burn whenever the \emph{osculating} SMA fell below the
lower edge, plus a trim burn. Two problems:
\begin{enumerate}
  \item \textbf{$J_2$ breaks the trigger.} The osculating SMA computed from the two-body energy
        $a_{\text{osc}} = -\mu/(2\varepsilon)$, $\varepsilon = v^2/2 - \mu/r$, swings by about 12\,km per
        orbit under $J_2$ (twice per orbit: high over the equator, low at the highest latitudes). A 500\,m
        band is therefore crossed every orbit: a controller watching the osculating SMA fires constantly.
  \item \textbf{``If'' statements are not ODEs.} An instantaneous impulse is a discontinuity that the
        solver handles only with event detection, and it does not represent a real (finite-thrust) burn.
\end{enumerate}

\subsection{The mean SMA}
Including the $J_2$ potential in the energy removes the swing:
\begin{equation}
  \varepsilon_{J_2} = \frac{v^2}{2} - \frac{\mu}{r} + U_{J_2}(\mathbf r),\qquad
  a_{\text{mean}} = -\frac{\mu}{2\,\varepsilon_{J_2}} .
\end{equation}
Since $J_2$ is conservative, $\varepsilon_{J_2}$ is constant under $J_2$ and changes only through drag and
thrust. To first order, $a_{\text{osc}} \approx a_{\text{mean}} - (J_2R_e^2/a)(3\sin^2\phi - 1)$: the swing
is $+6.4$\,km over the equator and $-5.8$\,km at $53^\circ$ latitude. The simulated swing matches this to
within 11\,m; the orbit-averaged osculating SMA sits a fixed $(J_2R_e^2/a)(1-\tfrac32\sin^2 i) = 275$\,m
above $a_{\text{mean}}$ (a definition, not an error).

\subsection{The smooth deadband law with memory}
The state gains a 7th component, the thruster on/off state $s\in[0,1]$. With the smooth switch
$\sigma(x) = \tfrac12\left(1+\tanh(x/w)\right)$ (width $w = 5$\,m of SMA),
\begin{equation}
  \tau\,\dot s = \underbrace{\sigma(a_L - a_{\text{mean}})(1-s)}_{\text{switch on below the lower edge}}
              - \underbrace{\sigma(a_{\text{mean}} - a_U)\,s}_{\text{switch off above the upper edge}}
              + \text{latch}(s),
\end{equation}
with a switch time $\tau=60$\,s and a weak latch term that pulls $s$ towards 0 or 1, so the state
\emph{remembers} whether it is burning: between the edges both switches are off and $s$ stays where it is.
This is a \textbf{hysteresis}: the same SMA (say mid-band) can be coasting or burning depending on the past,
which no memory-less function $s = f(a)$ can represent (step~9 measures exactly this). The thrust is
\begin{equation}
  \mathbf a_{\text{thrust}} = T\,\theta(s)\,\hat{\mathbf v},\qquad
  \theta(s) = \sigma\!\left(\frac{s - 0.9}{0.02}\right),\qquad T = 2\times10^{-4}\ \text{m/s}^2 ,
\end{equation}
so the throttle opens only once $s$ is past 0.9 (a Hall-effect-class thruster, $\sim$50\,mN on $\sim$250\,kg).

\figfile{step5_deadband/step5_summary.png}{0.98}{Step 5 summary: the reproduction of the brief's Figure 1
(mean SMA, controlled vs drag-only, and the miss distance), 400 days of station keeping, burns and fuel by
horizon, the hysteresis loop, the thruster ramp for several switch times, and fuel vs band width.}{fig:s5}

\subsection{Results over 400 days}
\begin{itemize}
  \item 102 burns, one every \textbf{3.905 days}, each \textbf{22.8\,min} long, $\Delta v = 0.273$\,m/s per burn
        (27.8\,m/s in total, 25.2\,m/s over 360 days).
  \item The mean SMA stays in $[6921.904,\,6922.402]$\,km for the whole run.
  \item The distance between the controlled orbit and its drag-only twin reaches 1800\,km after 15 days and
        saturates near the orbit diameter, because the two satellites drift to different phases of the
        same orbit (the brief's ``hundreds of km in a week'').
  \item Fuel hardly depends on the band width (3.8--4.2\,m/s over 60 days for 250--1000\,m bands): the fuel
        replaces what drag removes; a narrower band only means more, smaller burns.
\end{itemize}

% =====================================================================================================
\section{Step 0: bias--variance tuning of every setting}\label{sec:tuning}
% =====================================================================================================
Many methods below have a free setting (a stencil order, a smoothing window, a number of mixture
components, a sparsity threshold). Instead of picking them by hand, step~0 chooses each from data and
saves the choice in \code{.cache/tuned_settings_v1.json}; every later step reads it with
\code{settings.tuned(...)}, so the optimum is carried forward automatically.

\subsection{The decomposition}
Repeat an estimate $\hat y$ of a true value $y$ over many independent noise draws (Monte-Carlo seeds or
bootstrap re-fits):
\begin{equation}
  \underbrace{\mathbb E\big[(\hat y - y)^2\big]}_{\text{MSE}}
  = \underbrace{\big(\mathbb E[\hat y]-y\big)^2}_{\text{bias}^2}
  + \underbrace{\mathbb E\big[(\hat y-\mathbb E[\hat y])^2\big]}_{\text{variance}} .
\end{equation}
More flexible / less smoothed settings lower the bias and raise the variance; the best setting minimises
the MSE. For time series the decomposition is done sample by sample and averaged.

\subsection{The studies and what they chose}
\begin{center}\small
\begin{tabularx}{\textwidth}{llp{3.4cm}X}
\toprule
Setting & Default & Tuned & How\\\midrule
FD stencil order & 6 & per step: 1\,min$\to$4, 2\,min$\to$2, 5 and 10\,min$\to$6 & MSE of $c_d$ from 10\,m-noise positions at each sampling step\\
Savitzky--Golay window / order & 21 / 5 & 61 / 5 & MSE of the classifier's $\dot v$ (coast and burn weighted equally)\\
stacking window (step 8) & 21 / 5 & 41 / 5 & $\text{bias}^2 + (\text{noise})^2\,\text{variance}/N$: averaging $N$ burns removes variance, not bias\\
derivative method & SG & SG & walk-forward / walk-backward event F1 (below)\\
GMM components $K$ & 4 & 4 & bias$^2$ + variance of $P(\text{burn}\mid x)$ (Brier score), one-standard-error rule\\
SMA derivative window & 41 & 41 & MSE of the learned coast and burn \emph{rates}\\
STLSQ threshold & 0.05 & 0.05 & MSE of the SINDy law's predicted $\dot a$ on TEST (sparsest within 1\%)\\
\bottomrule
\end{tabularx}
\end{center}

\begin{whybox}
A single global stencil order was first chosen on 1-minute data; the figure audit then showed that the
10-minute Kalman comparison of step~4 had silently inherited it and become 10\% wrong. Truncation error
grows like $h^{p}$ while noise amplification falls with $h$, so the best order depends on the step $h$ ---
hence an order \emph{per sampling step}. The same logic gave step~8 its own (shorter) window: when 92 burns
are averaged the variance falls by 92 but the bias does not.
\end{whybox}

\subsection{Choosing the derivative by walk-forward and walk-backward testing}
For each way of differentiating the velocity data (Savitzky--Golay, central differences of order
2/4/6/8), a GMM is trained on a 40-day window and tested out of sample over the following (\emph{forward}:
train days 0--40) or preceding (\emph{backward}: train days 360--400) 15, 30, 60, 120, 240 and 360 days,
at 0.25, 0.5 and 1$\times$ the reference noise. At the reference noise only Savitzky--Golay keeps an event
F1 of 1.0 in all 12 periods (0 false events, 27\,s burn-start error); the central differences fail from
0.5$\times$ noise (8th order already at 0.25$\times$) because they amplify noise.

\figfile{step0_tuning/tuning_summary.png}{0.98}{Step 0 summary: bias$^2$, variance and MSE for each
setting, and the walk-forward / walk-backward winner.}{fig:s0}

% =====================================================================================================
\section{Step 2: learning the coefficients by regression}\label{sec:regression}
% =====================================================================================================
\subsection{Finite differences}
From positions $\mathbf r_k$ sampled every $h$, central differences of order $p$ give the velocity and
acceleration, e.g.\ for $p=2$:
\begin{equation}
  \dot{\mathbf r}_k \approx \frac{\mathbf r_{k+1}-\mathbf r_{k-1}}{2h},\qquad
  \ddot{\mathbf r}_k \approx \frac{\mathbf r_{k+1}-2\mathbf r_k+\mathbf r_{k-1}}{h^2},
\end{equation}
and wider stencils for $p=4,6,8$. The truncation error is $O(h^p)$; white position noise $\sigma$ becomes an
acceleration error of order $\sigma/h^2$ (amplified by the stencil weights). This is the central
bias--variance trade-off of the whole project.

\subsection{Linear least squares}
Each unknown multiplies a known ``shape'', so one linear solve recovers it:
\begin{align}
  \ddot{\mathbf r} &= \mu\,\Big(-\frac{\mathbf r}{r^3}\Big) &&\Rightarrow\ \hat\mu = (\Theta^\top\Theta)^{-1}\Theta^\top\ddot{\mathbf r},\\
  \ddot{\mathbf r} + \frac{\mu\mathbf r}{r^3} - \mathbf a_{J_2} &= c_d\,\big(-|\mathbf v|\mathbf v\big),\\
  \ddot{\mathbf r} + \frac{\mu\mathbf r}{r^3} &= [c_d,\ J_2,\ \mu_{\text{Moon}},\ \mu_{\text{Sun}},\ PC_RA/m]\cdot[\text{shapes}] .
\end{align}
The columns differ by up to 10 orders of magnitude (the Sun term is tiny), so they are scaled to unit length
before solving and the scaling is undone afterwards (same answer, no round-off trouble).

\figfile{step2_regression/step2_summary.png}{0.98}{Step 2 summary: $\mu$ vs sampling step (the Nyquist
limit $T/2 = 47.75$\,min), finite-difference error vs step and order, the correlation of the joint-fit
columns, and the joint fit vs data length (clean 1-min, 10\,m noise, clean 5-min).}{fig:s2}

\paragraph{Results.} On clean 1-min data the joint fit recovers $c_d$ and $J_2$ to $\sim$0.002\%. The Moon
($\sim$5\%) and especially the Sun (10--80\%) are poorly identifiable: their tidal shapes are small and
correlated with each other and with $J_2$ (column correlations 0.18--0.25). With 10\,m noise the errors
fall with data length (e.g.\ $c_d$ from 11\% at 15 days to 0.27\% at 360 days). On 5-minute data the
truncation error swamps the small terms (Moon and Sun 200--700\% off).

% =====================================================================================================
\section{Step 3: drag estimation and noise robustness}
% =====================================================================================================
\subsection{Three ways to get $c_d$}
\begin{enumerate}
  \item \textbf{FD regression} (above): two derivatives of noisy positions.
  \item \textbf{The energy method}: one derivative only. The mean SMA from the $J_2$ energy decays at
        \begin{equation}
          \dot a = -\frac{2a^2}{\mu}\,c_d\,|\mathbf v|^3 ,
        \end{equation}
        so $c_d$ follows from the \emph{slope} of a straight line fitted through the whole SMA history:
        $\hat c_d = -\dot a\,\mu/(2\bar a^2\,\overline{|v|^3})$. Fitting a line averages the noise of every sample.
  \item \textbf{Shooting}: propagate the orbit with trial parameters and fit the whole trajectory to the
        observed positions by Levenberg--Marquardt (started from the energy-method answer, because the cost
        is rugged far from the minimum).
\end{enumerate}

\begin{center}\small
\begin{tabular}{lcccccc}
\toprule
noise (15 days, 1-min) & 0\,m & 1\,m & 5\,m & 10\,m & 50\,m & 100\,m\\\midrule
FD regression $c_d$ error [\%] & 0.009 & 0.64 & 3.2 & 6.3 & 32 & 64\\
energy method $c_d$ error [\%] & 0.006 & 0.008 & 0.015 & 0.035 & 0.18 & 0.35\\
\bottomrule
\end{tabular}
\end{center}
With 10\,m noise the energy method improves from 0.03\% (15 days) to 0.001\% (360 days); shooting gives
0.05\% / 0.04\% at 15 / 30 days at a much higher cost.

\figfile{step3_drag_noise/step3_summary.png}{0.98}{Step 3 summary: noise study, $c_d$ vs data length, the
energy method's SMA trend, sensitivity maps (sampling step $\times$ noise) and the shooting cost landscape.}{fig:s3}

\subsection{Step 3b: realistic (coloured) noise}
Real tracking errors are not white. The realistic mixture combines white noise at four rates (every sample
and held for 10/100/1000 samples), pink ($1/f$) noise, a random walk ($1/f^2$) and once-per-orbit and
once-per-day sinusoids, scaled to the \emph{same RMS} as the white case so that only the spectrum differs.
Second derivatives multiply the noise spectrum by $f^4$ (fast noise hurts, slow noise is filtered); the
energy method fits a slow trend, so slow noise (drifts, random walk) can mimic drag decay. At 1\,m RMS over
30 days the realistic noise makes the energy-method $c_d$ 4.2$\times$ worse, SRP 4.0$\times$ and the Sun
2.3$\times$ worse, while FD $c_d$ and $J_2$ are almost unchanged.

\subsection{Workstream E: whole-trajectory cost surface}
The shooting fit of $(c_d, J_2)$ on 3 days of 10-min data has a cost surface (\code{inverse_cost_surface_3d})
with a narrow valley at the true $c_d$ running along $J_2$; the joint fit at 100\,m noise recovers $c_d$ to
0.006\% and $J_2$ to 0.0014\%, with a 27\,m RMS floor that is model error (the observed orbit also feels
the Moon, Sun and SRP, which the fitted model omits).

% =====================================================================================================
\section{Step 4: spectra, filtering and Kalman filtering}
% =====================================================================================================
\subsection{The periodogram and the $J_2$ signature}
The radial unmodelled acceleration (finite-difference acceleration minus point-mass gravity and drag) is
Fourier-transformed with a Hann window. $J_2$ depends on $\sin^2\phi$, so it appears at \textbf{twice the
orbital frequency}. The signal-to-noise ratio is the peak power within $\pm1\%$ of 2 cycles/orbit divided
by the median background power below 6 cycles/orbit (a fine sampling step must not be penalised by
high-frequency finite-difference noise far above the band of interest).
\begin{itemize}
  \item SNR falls with noise (15 days, 5-min: 47\,dB at 0.1\,km, 27\,dB at 1\,km, 12\,dB at 5\,km).
  \item SNR collapses once the sampling step exceeds the $J_2$ Nyquist limit $T/4 \approx 23.9$\,min
        (55\,dB at 1\,min, $-3$\,dB at 90\,min, step~4c).
  \item A longer record gives finer frequency resolution, but over months drag slowly shifts the line,
        so SNR peaks at $\sim$60 days.
\end{itemize}

\subsection{Filtering}
A band-pass around 2 cycles/orbit isolates $J_2$ and a low-pass below 0.3 cycles/orbit isolates drag. For
\emph{estimating} a force with a known shape, filtering does not beat least squares, which already is the
matched filter.

\subsection{The extended Kalman filter (EKF)}
The state $\mathbf x = (\mathbf r,\mathbf v, c_d)$ is propagated with the full gravity + $J_2$ + drag
model and its state-transition matrix $\Phi$ ($\dot\Phi = F\Phi$, $F$ the Jacobian); each noisy position
updates it:
\begin{align}
  \text{predict:}\quad & \mathbf x^- = f(\mathbf x),\quad P^- = \Phi P\Phi^\top + Q,\\
  \text{update:}\quad & K = P^-H^\top(HP^-H^\top+R)^{-1},\quad \mathbf x = \mathbf x^- + K(\mathbf z - H\mathbf x^-),\quad P = (I-KH)P^- .
\end{align}
With 10-min data and 10\,m noise the EKF finds $c_d$ to 0.04\% after 15 days and 0.01\% after 30 days,
better than the energy method (1.2\%) and FD regression (0.5--21\%) on the same coarse data.

\figfile{step4_spectral_kalman/step4_summary.png}{0.98}{Step 4 summary: where each force lives in frequency,
$J_2$ SNR vs data length and sampling step, band-pass and low-pass filtering, and the Kalman filter against
the other $c_d$ methods.}{fig:s4}

\subsection{Step 4b: ``force spectroscopy''}
Inspired by infrared spectroscopy, each force is identified by its \emph{diagnostic band} (a band where it
holds $\ge$90\% of the power) and its strength read from the band area relative to a reference library.
Parseval's theorem ($\sum|X_k|^2 \propto \overline{a^2}$) explains a key result: spectral \emph{power}
shares do not match the force budget (they square the amplitudes, so $J_2$'s 99.98\% becomes 99.999999\%);
the \emph{square root} of the power does (RMS $\approx$ mean $|a|$). Band reading is 30--40$\times$
faster than regression; it recovers drag to 0.5\%, $J_2$ to 2\% and SRP to 0.9\% in random mixtures, the
Moon to 8\%, while the Sun is not separable (its lines lie under the Moon's). A phase-aware fit is exact but
needs the force model evaluated along the orbit, like regression.

% =====================================================================================================
\section{Step 6: classifying coast and burn}\label{sec:classify}
% =====================================================================================================
\subsection{The feature: the unmodelled acceleration}
From measured $\mathbf r$, $\mathbf v$ (6-s samples with reference noise 0.1\,m and 0.5\,mm/s), the
unmodelled acceleration is the derivative minus the known physics:
\begin{equation}
  \mathbf a_u = \dot{\mathbf v}_{\text{measured}} - \Big(-\frac{\mu\mathbf r}{r^3} + \mathbf a_{J_2}\Big)
  \approx \mathbf a_{\text{drag}} + \mathbf a_{\text{thrust}} .
\end{equation}
While coasting $|\mathbf a_u|\sim 10^{-6}$\,m/s$^2$ (drag); during a burn $\sim 2\times10^{-4}$. Two features
are used: $x_1 = \log|\mathbf a_u|$ (how big) and $x_2 = \mathbf a_u\cdot\hat{\mathbf v}/10^{-4}$ (thrust
pushes \emph{forwards}, noise does not). Velocity-based ($\dot v$) features are used rather than
position-based ($\ddot r$) ones because a second derivative of noisy positions is pure noise.

\begin{whybox}
Why the derivative method matters: with the 2nd-order stencil at 6\,s the coast residual is the
stencil's own truncation error ($\sim6\times10^{-5}$\,m/s$^2$), only 4$\times$ below the thrust, so the GMM's
burn Gaussian was not centred on the burns. Higher-order stencils bring the floor to the real drag
($8\times10^{-7}$) on clean data; with noise only the Savitzky--Golay filter (a least-squares polynomial
through a window) works. The tuned window is 61 samples (6\,min).
\end{whybox}

\subsection{The Gaussian mixture model}
A GMM models the feature density as $p(\mathbf x) = \sum_{k=1}^{K}\pi_k\,\mathcal N(\mathbf x;\bm\mu_k,\Sigma_k)$,
fitted by expectation--maximisation on TRAIN. The posterior (``weight'') of a sample is
$P(k\mid\mathbf x) = \pi_k\mathcal N_k/\sum_j\pi_j\mathcal N_j$. With $K=4$ (tuned): one or two Gaussians for
the coast, one for the burn plateau and a spare for the thruster ramps; a component is labelled \emph{burn}
if its mean $\log|a_u|$ lies more than 3 coast standard deviations above the coast mean (thrust only adds
acceleration). Alternatives: a Bayesian GMM (Dirichlet prior switches off unused components), $k$-means,
Isolation Forest and a MAD threshold.

\subsection{From samples to events}
A 23-min burn is $\sim$230 samples. Flags are grouped into \textbf{events}: gaps of $\le2$ samples are
filled and runs shorter than 10 samples (1\,min) are dropped. Scoring is per event: a true burn is found if a
detected event overlaps it within $\pm60$\,s; precision = real events / detected events, recall = burns found
/ burns, F1 their harmonic mean.

\subsection{Results}
On 92 TEST burns over 360 days: GMM, Bayesian GMM and the MAD threshold score F1 = 1.0 on noisy data (all
92 burns, 0 false events), $k$-means fails completely with noise (136\,662 false events), Isolation Forest
0.85. The scores hold at every horizon from 15 to 360 days.

\subsection{Stress test: how coarse can the sampling get?}
At 0--4$\times$ the reference noise and steps from 6\,s to 60\,min:
\begin{itemize}
  \item the Savitzky--Golay derivative works up to a \textbf{2-min} step (beyond, a polynomial through 5
        samples spanning a large arc of the curved orbit has a truncation error bigger than the thrust);
  \item a new \textbf{one-step-propagation} feature --- propagate each measured state one step with gravity
        + $J_2$ (RK4) and divide the velocity mismatch at the next measurement by the step --- works from
        $\sim$1\,min up to \textbf{20\,min} at every noise level, since its noise shrinks as $1/h$;
  \item switching between the two at about 2\,min gives F1 $\ge$ 0.9 from 6\,s to 20\,min at all noise levels;
        beyond one burn length (23\,min) the burn is diluted into one interval.
\end{itemize}

\figfile{step6_classification/step6_summary.png}{0.98}{Step 6 summary: dataset split, $\ddot r$ vs $\dot v$
noise floors, the feature space with the GMM, the stress test, event F1 vs horizon, event scores, timing
errors and the working ranges of the sampling step.}{fig:s6}
\figfile{step6_classification/gmm_phase_space_3d.png}{0.7}{The phase space (mean SMA, its rate, along-track
acceleration) over two TEST cycles, every sample coloured by its GMM weight $P(\text{burn}\mid x)$.}{fig:phase}

% =====================================================================================================
\section{Step 7: robustness of the classifiers}
% =====================================================================================================
Days 0--40 train, the next 60 days test. GMM, Bayesian GMM and threshold keep F1 = 1.0 from 0.25 to
8$\times$ the reference noise; $k$-means fails as soon as there is noise (its two clusters split the
coast cloud); Isolation Forest stays near 0.8--0.9. The GMM needs $\sim$1000 training samples. The
number of components hardly matters between 2 and 6 at these noise levels. The workstream \code{detection}
stage adds five basic detectors (rolling z-score, MAD threshold, $k$-means, GMM, CUSUM) on the same data:
CUSUM is the most noise-tolerant online detector.

% =====================================================================================================
\section{Step 8: repetition --- autocorrelation, stacking and a Gaussian process}
% =====================================================================================================
\begin{itemize}
  \item The \textbf{autocorrelation} of the noisy along-track acceleration peaks at 3.9051 days, matching the
        true burn period (3.9051 days).
  \item \textbf{Stacking}: aligning the 92 detected burns and averaging them reduces noise; the error of the
        average falls until it hits a floor set by the smoothing of the edges (the tuned stacking window
        trades this bias against the $1/\sqrt N$ variance).
  \item A \textbf{Gaussian process} (RBF + white-noise kernel) through the stacked samples gives a thrust law:
        level $2.00\times10^{-4}$\,m/s$^2$ (0.999 of the truth), duration 1368\,s (exact), rise time 60\,s
        (truth 30\,s: the derivative smooths the edge).
  \item With only 15 days (4 burns) the level is already within 0.3\%, the period exact to $10^{-4}$.
\end{itemize}
\figfile{step8_repetition_gp/step8_summary.png}{0.98}{Step 8 summary.}{fig:s8}

% =====================================================================================================
\section{Step 9: learning the control law (SINDy and BINDy) and forecasting}\label{sec:sindy}
% =====================================================================================================
\subsection{Sparse identification of nonlinear dynamics (SINDy)}
SINDy writes the rate of a state as a sparse combination of candidate functions (a ``library''):
\begin{equation}
  \dot a = \Theta(a,s)\,\bm\xi,\qquad \Theta = [\,\theta_1(a,s),\ \theta_2(a,s),\ \dots\,],
\end{equation}
and finds a sparse $\bm\xi$ by \textbf{sequentially thresholded least squares} (STLSQ): solve least squares,
set coefficients below a threshold to zero, refit on the remaining terms, repeat.

\paragraph{Why a cluster-indicator library.} The obvious physics library $[D(a), B(a), 1]$ with
$D = \tfrac{2a^2}{\mu}|v|^3$ (drag shape) and $B = \tfrac{2a^2}{\mu}|v|$ (thrust shape) is singular: over a
500\,m band both are constant to a few ppm and perfectly correlated (correlation 0.997). The library
therefore uses the thruster state from the classifier: $[1-s,\ s,\ a-a_{\text{ref}},\ s(a-a_{\text{ref}}),\
(a-a_{\text{ref}})^2]$. STLSQ keeps only the two cluster terms; drag and thrust follow from them:
$\hat c_d = 1.0005\,c_d$ and $\hat T = 0.9997\,T$.

\paragraph{BINDy (Bayesian).} Automatic relevance determination (ARD) regression gives a posterior mean and
standard deviation for each coefficient, i.e.\ credible intervals for the physics.

\subsection{The switching law}
The burn start and end SMAs give the learned edges 6921.9037 / 6922.4013\,km. A memory-less sigmoid
$s = f(a)$ reaches only 0.50 balanced accuracy (chance); the hysteresis with the learned edges reaches
1.00. \textbf{The controller has memory, and the learned law must too.}

\subsection{Forecasting}
The learned hybrid model --- $\dot a$ from SINDy, switching by the learned hysteresis --- is integrated
forward from day 40 with no further data. Mean onset error over the 360-day forecast: periodic baseline
0.9\,min, GP coast curve 18\,min, BINDy ensemble 34\,min, SINDy 195\,min, random forest 286\,min.

\begin{whybox}
Why does SINDy drift and why does the periodic baseline win here? A coast rate learned from 40 days of
noisy $\dot a$ carries a few hundred ppm of error; each 3.9-day coast is then a little too short, and the
error accumulates linearly over 92 cycles. The periodic baseline learns the period directly, which is
perfect \emph{only because drag is constant in this simulation}. Step~10 removes this assumption and the
ranking reverses.
\end{whybox}
\figfile{step9_sindy_control_law/step9_summary.png}{0.98}{Step 9 summary.}{fig:s9}

% =====================================================================================================
\section{The workstream stages, aligned with the report data}
% =====================================================================================================
The Week~2--5 stages were rewritten to use the report's data, split, features and event scoring, and
gained 3D views:
\begin{description}
  \item[\code{neural_net}] an MLP fitted as SMA$=f(t)$ on days 0--40 flattens outside its training window
        (44\,km off after 360 days) whereas a straight line is 0.09\,km off: a time-indexed network
        cannot extrapolate physics.
  \item[\code{inverse}] the 3D cost surface over $(c_d, J_2)$ (above).
  \item[\code{spectral}] the $J_2$ SNR surface over sampling step $\times$ noise and a waterfall of the
        spectrum against record length.
  \item[\code{smooth_deadband}, \code{deadband}] Figure~1, the switches, the $J_2$ mean-vs-osculating figure, the
        400-day sawtooth and a 3D hysteresis loop.
  \item[\code{derivative_gmm}] the GMM in $\ddot r$/$\dot v$ space, five derivative methods, the Week~4
        features on the same samples, and a cadence sweep.
  \item[\code{gmm}] the GMM density and burn-probability surfaces in 3D, the phase space coloured by GMM
        weight, and GMM vs $k$-means decision regions.
  \item[\code{heldout}] seven detectors on held-out days 40--100, including a Gaussian-process anomaly
        detector that models the coast SMA as a function of the time since the SMA left the upper edge
        (trained on GMM labels, not truth): GMM, Bayesian GMM, DBSCAN and the GP reach F1 = 1.0,
        Isolation Forest 0.71, $k$-means and Ward clustering 0.
\end{description}

% =====================================================================================================
\section{Extensions towards the brief (steps 10--18)}
% =====================================================================================================
\subsection{The mean-element model}
A full 6-s integration of 400 days takes about 20 minutes; the extensions need hundreds of scenarios. They
use the model the controller itself follows:
\begin{equation}
  \dot a = -k\,w(t)\,e^{-(a-a_{\text{ref}})/H} + T\,s,\qquad s:\ 0\to1 \text{ at } a\le a_L,\ 1\to0 \text{ at } a\ge a_U ,
\end{equation}
with $k$ the coast decay rate (5.335\,m/h), $T$ the thrust contribution (1308\,m/h), $H=60$\,km the
density scale height and $w(t)$ the space-weather density factor. $k$, $T$ and the \emph{effective} edges
(the tanh switch starts each burn $\sim$6\,m above the nominal lower edge) are calibrated on the full run;
the model reproduces the burn period to 0.01\% and the burn length to 1\%. The effective band is the SMA span covered by a coast (coast rate $\times$ coast duration), about 3\,m wider than the SMA difference between the throttle half-way points because of the thrust ramps.

\subsection{Step 10: space-weather drag}
$w(t)$ combines a 27-day solar-rotation wave, a slow solar-cycle trend and 8 geomagnetic storms (density
$\times$2--4, 1.5-day decay); an observable daily proxy (an F10.7/Ap-like index) has 5\% noise. Forecasters:
periodic baseline; the law with constant drag; with the latest coast rate; with $k\times$proxy and a 27-day
persistence forecast of the proxy; and with a perfect space-weather forecast (upper bound). Forecasts are
event-driven: with $G(t)=\int k\,\hat w\,dt$ a coast from $a_0$ ends where $G(t)-G(t_0)=a_0-a_L$.

\begin{center}\small
\begin{tabular}{lcccc}
\toprule
mean $|$onset error$|$ [min] & const., 15\,d & const., 360\,d & space weather, 15\,d & space weather, 120\,d\\\midrule
periodic baseline & 2 & 3 & 1642 & 1732\\
law, constant drag & 68 & 630 & 1418 & 1670\\
law + space-weather proxy & 68 & 630 & 730 & 1286\\
law + perfect forecast & 68 & 630 & 164 & 419\\
\bottomrule
\end{tabular}
\end{center}
\textbf{Variable drag changes the interval between burns but not the band edges}: the periodic baseline then
loses its phase, while the learned law tracks it as well as drag can be forecast.
\figfile{step10_space_weather/step10_summary.png}{0.98}{Step 10 summary.}{fig:s10}

\subsection{Step 11: realistic tracking passes}
Instead of 6-s states, 8-minute passes of 12-s positions (10\,m--1\,km noise), 1--8 per day. Each pass is
fitted by batched Gauss--Newton orbit determination (6 state parameters, RK4 gravity + $J_2$, finite-difference
Jacobian); one pass gives the SMA to about 2$\times$ the position noise. Burns are found in the gaps:
the SMA jump minus the drag decay (estimated from pass pairs a day apart) must exceed a threshold set by
the pass noise. Pass-to-pass: F1 $\approx$1 at 10\,m noise for 1--8 passes/day; at 100\,m only a multi-pass
test (mean of all passes a day before vs after) finds some burns (F1 0.47 at 8/day); at $\ge$300\,m the
500\,m burn is buried. A CCSDS OEM reader runs the same analysis on real ephemerides placed in
\code{data/ephemeris/}.

\subsection{Step 12: close-approach screening}
An SMA error $\delta a$ changes the mean motion by $\delta n = -\tfrac32 (n/a)\,\delta a$, so the along-track
error grows as
\begin{equation}
  \delta s(t) = -\tfrac32\,n\int_0^t \big(a_{\text{pred}}-a_{\text{true}}\big)\,dt ;
\end{equation}
an unknown 500\,m burn makes the catalogue drift by $\tfrac32 n\times0.5$\,km $\approx$ 3\,km/hour. RMS
along-track error at 1 / 7 days: drag-only 19 / 396\,km, periodic 18 / 42\,km, learned law with
space weather 5.8 / 11.9\,km, perfect forecast 2.0 / 4.1\,km. Screening with $P_c \approx \pi R^2\,
\mathcal N(\text{miss};0,\Sigma)$ ($R = 20$\,m, alert at $10^{-4}$): with an honest covariance even the best
method misses 25\% of dangerous conjunctions at 1 day (drag-only 68\%), because the errors are still km-sized
--- the ``dilution'' that the brief warns makes screening ineffective. With today's catalogue-only covariance
nearly all dangerous conjunctions are missed.

\subsection{Step 13: sensor tasking and track association}
Six observation opportunities a day, a budget of 1, 2 or 4. With an association gate of 20\,km, custody is
lost after 74\% of burns with one random observation a day, 48\% if the observation is scheduled just after
each \emph{predicted} burn, and 15--19\% if the catalogue also carries the predicted burn (its error is then only
$\tfrac32 n\,\delta a\times$ timing error). With 4 observations a day the losses are 20\% and 4\%.

\subsection{Step 14: several control laws on one satellite}
Orbit raising from 6890\,km, the 500\,m deadband, a 3\,km band change at day 150 and six 3-min collision-avoidance
burns. From noisy SMA (5\,m, 10-min samples) every burn is detected and labelled from duration, SMA gain and
start SMA relative to a running edge: \textbf{105/105 correct} by rules and also by an unsupervised GMM; the
band change is found as a changepoint at day 150.0; each regime gets its own rate law.

\subsection{Step 15: a fleet of 40 satellites}
Forty satellites share the space weather but differ in band ($\pm$5\,km), drag ($\times$0.7--1.3) and thrust
($\pm$10\%). Three pattern-of-life changes are planted (band up 1\,km at day 200, thrust halved at day 250,
band narrowed to 200\,m at day 300). Rolling 30-day laws are compared with each satellite's own history,
with the coast rate divided by the fleet median (removing the common space weather): exactly the three
changed satellites are flagged, 10--20 days after each change, with no false flags. Empirical-Bayes pooling of
the drag coefficient, $\hat k_i^{\text{pool}} = (\tau^2\hat k_i + s_i^2\bar k)/(\tau^2+s_i^2)$, reduces the
error from 21.6\% to 12.5\% with 4 days and from 11.8\% to 6.8\% with 8 days of data per satellite
(with only 2 days the single-arc variance has to be guessed and pooling overshoots).

\subsection{Step 16: calibrated (conformal) burn-time intervals}
Split conformal prediction: the absolute onset errors of forecasts from calibration origins (days 40--200)
give, per lead-time bucket, the $\lceil (n+1)(1-\alpha)\rceil/n$ quantile $q$; the interval is $\pm q$.
On new origins (days 200--360) the 90\% intervals of the space-weather forecaster cover 93--99\% of the burns.

\subsection{Step 17: GEO station keeping}
The GEO longitude obeys $\ddot\lambda = -A\sin 2(\lambda-\lambda_s)$ (Earth's triaxiality, $A=0.0017^\circ$/day$^2$,
$\lambda_s = 75.07^\circ$E); the inclination grows by 0.85$^\circ$/year. East-west (impulsive, $\pm0.05^\circ$ box),
north-south (impulsive) and an electric-propulsion east-west law (continuous low thrust with a deadband on
position and drift rate) are simulated. SINDy recovers $A = 0.00171$ and $\lambda_s = 74.98^\circ$ from a
free-drift phase; inside the box only the local acceleration is identifiable. Impulses are detected from
6-hourly longitudes as drift-rate jumps between 8-day quadratic fits, and the east-west burns are forecast
0.1--1.3 days off up to 225 days ahead from the learned local acceleration and box edges.

\subsection{Step 18: the SINDy equation zoo}
Ten candidate libraries per dataset --- constant per cluster, polynomial, exponential density, time
polynomial, Fourier in time, space-weather proxy, proxy $\times$ density, a deliberately wrong library without
a cluster indicator and a ``kitchen sink''; for GEO polynomials and Fourier harmonics in $\lambda$:
\begin{itemize}
  \item on the full-physics data every richer library is pruned back to the same two-term law
        (the data support nothing more);
  \item with space weather only the proxy libraries forecast well (about 260 vs $\sim$1640\,min);
  \item for GEO the true $\sin2\lambda,\cos2\lambda$ form (with or without a constant) forecasts best;
  \item ``kitchen-sink'' libraries diverge through collinearity; BIC alone can prefer an over-fitted library
        that forecasts worse, so free-running forecasts are the decisive test;
  \item \textbf{thresholding}: classic STLSQ drops terms small \emph{relative to the largest} coefficient,
        which deletes the real coast drag next to the 250$\times$ larger, rare burn term; a significance
        threshold ($|\xi|\ge 3$ standard errors) keeps it but lets spurious terms through on very large,
        correlated data sets. Both are reported side by side.
\end{itemize}
\figfile{step18_sindy_zoo/step18_summary.png}{0.98}{Step 18 summary.}{fig:s18}

% =====================================================================================================
\section{Step 19: thrust law and natural dynamics, learned together or apart}\label{sec:split}
% =====================================================================================================
\subsection{Why split the problem?}
Step~9 learned everything at once. But the mean-SMA dynamics are the sum of two physically different parts,
\begin{equation}
  \dot a \;=\; \underbrace{-\,c_d\,D(a,|\mathbf v|)}_{\text{natural (coasting) dynamics}}
         \;+\; \underbrace{T\,\theta(t)\,B(a,|\mathbf v|)}_{\text{thrust control law}},\qquad
  D = \frac{2a^2}{\mu}|\mathbf v|^3,\quad B = \frac{2a^2}{\mu}|\mathbf v|,
\end{equation}
plus the switching law that decides $\theta$ (on at the lower edge, off at the upper edge). In
operations one part is often better known than the other: the natural dynamics come from a catalogue
force model, while the operator's thruster and control law are unknown --- or, for one's own satellite,
the thrust is known and the drag is not. Step~19 measures what each kind of prior knowledge is worth.

\subsection{The four scenarios}
All use step~9's data (the 400-day controlled run, reference noise, smoothed derivative $\dot a$,
burn flags from the GMM), learn on days 0--40 and are scored on days 40--400 against the \emph{true}
model ($c_d$, $T$, the true throttle profile and the true effective edges). Each is fitted by least
squares (SINDy point estimate) and by ARD regression (BINDy posterior standard deviation).
\begin{description}
  \item[A. Whole system] one regression of $\dot a$ on $[-D,\ sB]$ over \emph{all} samples, with $s$ the
        detected burn flags; edges from the detected burns. Everything learned at once, as it comes.
  \item[B. Thrust law only (natural known)] $c_d$ is given: $\dot a + c_dD = T\,sB$ gives $T$; edges learned.
  \item[C. Natural law only (thrust known)] $T$, the throttle profile and the edges are given:
        $\dot a - T\theta B = -c_dD$ over all samples (burns included) gives $c_d$.
  \item[D. Both, learned separately] $c_d$ from the clean \emph{coast} arcs only, then $T$ from the clean
        \emph{burn} arcs given that $c_d$.
\end{description}
The forecast integrates the hysteresis law forward from the end of the last training burn with the
learned rates (shapes averaged over the training coast and burn samples, i.e.\ with the measured, not
circular, speeds) and compares the predicted burn onsets with the true ones over 360 days.

\subsection{Results}
\begin{center}\small
\begin{tabular}{lcccc}
\toprule
scenario & $c_d$/true (BINDy $\pm1\sigma$) & $T$/true & period [d] & onset error at 360\,d\\\midrule
A whole system & 0.9922 ($\pm$1.4\%) & 0.924 & 3.9351 & 1668 min\\
B thrust only (natural known) & given & 0.924 & 3.9047 & 26 min\\
C natural only (thrust known) & 1.0015 ($\pm$1.0\%) & given & 3.9006 & 302 min\\
D both, separately & 1.0006 ($\pm$0.6\%) & 0.9995 & 3.9011 & 262 min\\
\midrule
truth & 1 & 1 & 3.9051 & ---\\
\bottomrule
\end{tabular}
\end{center}
\begin{itemize}
  \item \textbf{The joint regression is biased.} The GMM's burn flags are slightly wider than the
        thrust (they include the ramps, and the smoothed derivative spreads the burn edge over the
        filter window), so in scenario A some samples with partial thrust are labelled ``burn'' and
        some with thrust are labelled ``coast''. Least squares then lowers $T$ by 7.6\% and absorbs part
        of the edge leakage into the drag ($-0.8\%$). Splitting the problem onto clean arcs (D) removes
        both biases: $c_d$ to 0.06\%, $T$ to 0.05\%.
  \item \textbf{The forecast is governed by the natural law.} A coast is 99.6\% of each 3.9-day cycle,
        so a 0.1\% drag error shifts every burn by about 5 minutes and accumulates linearly. With the
        natural law known (B) the forecast stays within 26 minutes over a year, even though its $T$ is
        7.6\% low (that only lengthens each 23-minute burn by about 2 minutes). With the thrust known (C)
        the drag is learned to 0.15\% and the forecast is far better than A but worse than B.
  \item \textbf{Rate errors.} Predicted $\dot a$ on coasting samples is within 0.001--0.04\,m/hour of
        the truth for all scenarios; over all samples the error (about 20\,m/hour) is dominated by the
        few samples at the burn edges where the detected flags and the true throttle disagree --- except
        in C, which uses the true throttle (0.008\,m/hour).
\end{itemize}
\begin{keybox}{What prior knowledge is worth}
For forecasting a station-keeping satellite, knowing the \emph{natural} dynamics (a good drag model)
is worth far more than knowing its thruster: the thrust level hardly matters, the coast rate decides
when the next burn comes. When nothing is known, learn the two laws \emph{separately} on clean arcs
rather than in one regression.
\end{keybox}
\figfile{step19_sindy_split/step19_summary.png}{0.98}{Step 19: (a) physics recovered per scenario with
BINDy intervals, (b) free-running forecast error, (c) one burn as seen by each scenario, (d) the free
decay without a deadband, (e) the learned decay laws, (f) their SMA forecasts.}{fig:s19}

\subsection{The natural law with no deadband}
Without the controller the satellite decays freely: 51\,km in 400 days for both the drag-only run
($J_2$ + drag) and the natural run (plus Moon, Sun and SRP). With constant density the true law is
\begin{equation}
  \dot a = -c_d\,\frac{2a^2}{\mu}|\mathbf v|^3 = -2\,c_d\sqrt{\mu\,a}\qquad\text{(circular orbit)} .
\end{equation}
Eight libraries were fitted on days 0--40 (one SMA sample per minute, reference noise, 1-hour smoothing
derivative) and tested on days 40--400:
\begin{center}\small
\begin{tabular}{lccc}
\toprule
library & test RMS of $\dot a$ [m/h] & SMA error after 360 d & recovered $c_d$\\\midrule
constant / linear / quadratic$^{*}$ & 0.012 & 84 m & ---\\
$\sqrt{a}$ (true form) & 0.0006 & 1 m & 1.0001\\
physics shape $D(a)$ (true form) & 0.0006 & 1 m & 1.0001\\
exponential density $e^{-(a-a_0)/H}$ & 3.55 & 51 km & ---\\
power laws $a^{-1}, a^{1/2}, a^2$ & 0.14 & 0.9 km & ---\\
kitchen sink & 0.14 & 0.9 km & ---\\
\bottomrule
\end{tabular}\\[2pt]
{\footnotesize $^{*}$ STLSQ prunes the linear and quadratic terms, leaving a constant rate.}
\end{center}
\begin{itemize}
  \item The \textbf{true form} recovers $c_d$ to 0.01\% from 40 days and forecasts the SMA a year ahead to
        1\,m (also with the Moon, Sun and SRP present: their effect on the mean SMA is periodic).
  \item A \textbf{constant} rate is almost as good over 40 days but misses the 0.4\% change of the rate over
        the 51\,km descent: 84\,m after a year. STLSQ correctly prunes the linear and quadratic terms ---
        over 40 days they cannot be distinguished from noise.
  \item The \textbf{exponential-density} library encodes physics that is \emph{not} in this simulator
        (density rising as the orbit falls; here density is constant): it fits the training data but its
        extrapolation is 51\,km off --- a physically plausible library is not automatically the right
        one, and free-running forecasts, not training fit, decide.
  \item Collinear \textbf{power-law} and kitchen-sink libraries split the rate between near-identical
        columns and extrapolate to 0.9\,km error.
\end{itemize}
In the real atmosphere the density \emph{does} rise as the orbit decays (and varies with space weather,
step~10), so on real data the exponential library would be expected to win; the method, not the winner,
is the transferable result.

% =====================================================================================================
\section{Validation, figure audit and honest limitations}
% =====================================================================================================
\subsection{Testing}
41 unit tests (\code{python -m unittest discover tests}) cover the physics ($J_2 = -\nabla U_{J_2}$), the
controller, finite differences, the bias--variance identity, the mean-element calibration, forecasting with
perfect knowledge, conformal coverage, $P_c$ against Monte Carlo, pass orbit determination, the OEM reader,
regimes, fleet change detection, GEO law recovery and the SINDy zoo.

\subsection{Problems found and fixed during the audits}
\begin{itemize}
  \item the $J_2$ reference radius (0.22\% too weak);
  \item a $J_2$ check that ignored the drag-driven change of $a$ (1.4\% apparent error);
  \item a GMM burn component not centred on the burns (2nd-order derivative truncation error);
  \item a single tuned stencil order applied to 10-min data (10\% $c_d$ error, now per step);
  \item an over-smoothed stacked burn in step~8 (now its own tuned window);
  \item burn edges biased by partially-thrusting samples (now from coast/burn line intersections);
  \item a GP detector that used true labels (now GMM labels);
  \item integrators overshooting the band edges (now stopped on the edge).
\end{itemize}

\subsection{Limitations}
\begin{itemize}
  \item The inputs of steps 1--9 are full state vectors at 6\,s with 0.1\,m / 0.5\,mm/s noise --- far better
        than real tracking (step~11 shows how quickly detection degrades with realistic passes).
  \item Constant drag in steps 1--9 makes the periodic baseline look unbeatable; step~10 corrects this
        with modelled, not real, space weather.
  \item Steps 10--16 use the mean-element model (burn period within 0.01\% of the full run), not full orbit integration.
  \item No real ephemerides were downloaded (a space-track account is needed); optical angles-only orbit
        determination is not implemented.
  \item The per-step stencil order was tuned for $c_d$ under noise; at the 2-min step it picked order 2, which
        is worse for $\mu$ on clean data (step~2a) --- tuning is objective-specific.
  \item Moon, Sun and SRP strengths are only weakly identifiable from short or coarse records.
\end{itemize}

% =====================================================================================================
\section{How well the project answers the brief}
% =====================================================================================================
\begin{center}\small
\begin{tabularx}{\textwidth}{lX}
\toprule
Brief requirement & Status\\\midrule
Figure 1 & reproduced (step 5) \\
Detect past manoeuvres & 92/92 burns, 0 false events over 360 days at reference noise; down to 10\,m-noise tracking passes\\
Reconstruct the control law & edges to $\sim$0.3\,m, drag and thrust to $<$0.05\%, thrust profile, memory (hysteresis), multiple regimes, GEO law\\
Predict future manoeuvres & yes; beats a periodic baseline once drag varies, with calibrated intervals\\
Screening / association / tasking & demonstrated quantitatively (steps 12--13) in a simplified setting\\
Real data, realistic noise & partly: tracking passes and an OEM reader; no real ephemerides yet\\
\bottomrule
\end{tabularx}
\end{center}

\appendix
\section{Running the code}
\begin{verbatim}
python main.py                    # all 31 stages -> outputs/ (~30 min)
python main.py --list             # the stages, in run order
python main.py --stages step6 gmm # only some stages
python main.py --retune           # redo the step 0 bias-variance studies
python -m unittest discover tests # 41 tests
\end{verbatim}
Library code is in \code{deadband/} (no plotting); each stage in \code{stages/} draws figures into
\code{outputs/report/stepN_*/} and returns numbers for \code{outputs/results.json}. Every module starts with
a plain-English docstring explaining its method.

\section{Glossary}
\begin{description}[style=nextline]
  \item[SMA] semi-major axis; \emph{osculating} from the two-body energy, \emph{mean} including the $J_2$ potential.
  \item[Deadband] the allowed SMA band; the controller burns only at its edges.
  \item[Hysteresis] dependence of the state on its history (the controller's memory).
  \item[GMM] Gaussian mixture model; EM = expectation--maximisation.
  \item[Event F1] harmonic mean of event precision and recall.
  \item[SINDy / STLSQ] sparse identification of nonlinear dynamics / sequentially thresholded least squares.
  \item[BINDy / ARD] the Bayesian variant / automatic relevance determination.
  \item[EKF] extended Kalman filter; STM = state-transition matrix.
  \item[TCA, $P_c$] time of closest approach, collision probability.
  \item[Conformal prediction] distribution-free intervals calibrated on held-out errors.
  \item[OEM] CCSDS orbit ephemeris message (ephemeris file format).
\end{description}

\end{document}
